\documentclass[10pt,a4paper]{amsart}
\usepackage[margin=19mm]{geometry}
\usepackage{amsmath,amssymb,mathtools}
\usepackage[scr=rsfs,cal=euler]{mathalfa}
\usepackage{xfrac}
\usepackage[unicode,hidelinks,bookmarks=false]{hyperref}
\newcommand{\ie}{i.e.\ }
\newcommand{\cf}{cf.\ }
\newcommand{\resp}{respectively\ }
\newcommand{\Subsection}{\subsection}
\providecommand{\nobreakdash}{\nobreak\hbox{-}}
\setlength{\emergencystretch}{2em}
\multlinegap0pt
\usepackage{bm}
\usepackage{bbm}
\usepackage{dsfont}
\usepackage{xspace}
\usepackage{cancel}
\usepackage{mathtools}
\usepackage{amsthm}
\makeatletter
\@ifundefined{theorem}{\newtheorem{theorem}{Theorem}}{}
\@ifundefined{lem}{\newtheorem{lem}[theorem]{Lemma}}{}
\makeatother
\title[Robust Patrol in a Dispersed Environment]{Robust Patrol in a Dispersed Environment}
\author{Edward Mellor, Kevin Glazebrook, Kyle Lin, Rob Shone}
\date{Source: arXiv:2608.10935v1 (CC BY 4.0)}
\begin{document}
\maketitle

\noindent\emph{Source and licence.} Robust Patrol in a Dispersed Environment. Version arXiv:2608.10935v1 (CC BY 4.0), distributed under the Creative Commons Attribution 4.0 International licence, as recorded in the arXiv licence metadata for this exact version and shown on the abstract page https://arxiv.org/abs/2608.10935v1. Full rights notice: licenses/arxiv-2608.10935-CC-BY-4.0.txt.\par\smallskip

\noindent\textit{Editorial scope.} Counted unit: the Theorem whose source label is \texttt{th:one\_location} in the article (source0.tex lines 283--288, source label \texttt{th:one\_location}), delivered together with the article's own complete original proof of it (source0.tex lines 290--347, one \texttt{proof} environment of 58 lines). No mathematics is written, repaired or re-derived: statement and proof are the article's own text line for line. Required context retained uncounted: lem: am>=hm (lines 256--264). Every definition, display and label of those lines is retained unchanged. Omitted from this package and not redistributed: 1--255, 265--282, 289--289, 348--1394. Nothing inside a retained excerpt is omitted or altered: every retained body is delivered exactly as the article's lines are written, so no omission receipt of any kind accompanies this package. Citations: the retained excerpts carry no citation command at all, so no bibliography entry of the article is transcribed and no reference is redistributed. Reference adaptation: none was needed and none was made; every reference inside the delivered unit resolves inside it, so no unresolved reference or citation is produced. Printed slips kept literal: no printed word, symbol, hypothesis or number of the counted statement or of its proof was altered. Layout and class: the article's own class is \texttt{article} with packages including inputenc, amsmath, amsfonts, amssymb, amsthm, hyperref, graphicx, caption, subcaption, float, mathtools, url, algorithm, algorithmic; the wrapper is the lane's pinned \texttt{amsart} editorial wrapper at 10pt on A4, which restates the theorem-like environments the retained text needs and adds the article's own macro definitions that the retained text needs (none), with extra wrapper packages (bm, bbm, dsfont, xspace, cancel, mathtools). The delivered numbering is the unit's own. Assets: no retained span contains a figure environment, a graphics-inclusion command, a PDF-page inclusion, a table or a TikZ picture; no image file is copied into this package, the package declares no asset dependency and no image file is redistributed with it. Licence: version arXiv:2608.10935v1 of this article is distributed under the Creative Commons Attribution 4.0 International licence, as recorded in the arXiv licence metadata for this exact version and shown on the abstract page https://arxiv.org/abs/2608.10935v1. A full rights notice is delivered at licenses/arxiv-2608.10935-CC-BY-4.0.txt. Characters: every retained excerpt is pure ASCII, so the delivered engine, XeLaTeX with its default Latin Modern text font, reports no missing character.
\begin{lem}
\label{lem: am>=hm}
Suppose $g(x)$ is a function defined on $x \in [0,c]$ for some $c \in \mathbb{R}^+$.
If there exist $a,b \in \mathbb{R}^+$ with $a < b$ such that $g(x) \in [a,b]$ for $x \in [0,c]$, then
\[
\frac{\int_0^c g(x) dx }{c} \geq \frac{c}{\int_0^c \frac{1}{g(x)} dx},
\]
with equality if and only if $g(x)$ is constant for $x \in [0,c]$.
\end{lem}

\begin{theorem}
\label{th:one_location}
Consider a particular location and write $\lambda$ for the patroller's instantaneous detection rate at this location.
If the long-run fraction of effort the patroller spends at this location is capped at $r \in (0,1)$, then to minimize the expected time to detect an attacker at this location, it is optimal for the patroller to continuously allocate the same fraction of effort $r$ at this location at all times.
The minimized expected time to detection is $1 / (r \lambda)$.
\end{theorem}

\begin{proof}
An arbitrary policy can be delineated by some $c > 0$ and a function $h(t) \in [0,1]$ defined for $t \in [0,c]$ with the interpretation that the patroller allocates $h(t) \in [0,1]$ fraction of their effort at time $t \in [0,c]$ such that
\begin{equation}
\frac{\int_0^c h(t) dt}{c} = r,
\label{eq:r}
\end{equation}
and repeats the same pattern of length $c$ indefinitely.
It follows from renewal reward theory that the long-run fraction of effort allocated to this location is $r$.


Because an attacker does not know the patroller's real-time location and chooses the arrival time arbitrarily, the attacker is equally likely to arrive at any point during the cycle.
Write $g(t)$ for the expected time for the patroller to detect an attacker that arrives at time $t \in [0,c]$, so we seek to show that choosing $h(t) = r$ minimizes
\begin{equation}
\int_0^c g(t) \; \frac{1}{c} \; dt = \frac{\int_0^c g(t) dt}{c}.
\label{eq:obj}
\end{equation}
Consider an attacker arriving at time $t$ and condition on whether the attacker is detected in the interval $[t, t+ \delta)$ to get
\[
g(t) = \delta  + (1 - \lambda h(t) \delta + o (\delta)) \;  g(t+\delta),
\]
because with probability $1 - \lambda h(t) \delta + o (\delta)$ the attacker is still at the location at time $t + \delta$ so the expected additional time to detection is $g (t + \delta)$.
Rearrange the preceding to get
\[
\frac{g(t+ \delta) - g(t)}{\delta} = \left( \lambda h(t) + \frac{o(\delta)}{\delta} \right) g(t+ \delta) - 1.
\]
Taking $\delta \rightarrow 0$ yields
\[
g'(t) = \lambda h(t) g(t) -1.
\]
Solve the preceding for $h(t)$ to get
\[
h(t) = \frac{g'(t)}{\lambda g(t)} + \frac{1}{\lambda g(t)}.
\]
Because the allocation function $h(t)$ needs to meet the constraint in \eqref{eq:r}, we have that
\begin{align*}
r c 
&= \int_0^c h(t) dt \\
&= \int_0^c \frac{g'(t)}{\lambda g(t)} dt +  \int_0^c  \frac{1}{\lambda g(t)} dt \\
&= \frac{1}{\lambda} \ln g(t) \big|_0^c + \frac{1}{\lambda} \int_0^c \frac{1}{g(t)} dt \\
&= 0 + \frac{1}{\lambda} \int_0^c \frac{1}{g(t)} dt.
\end{align*}

The last equality follows because $g(c) = g(0) \geq 1/\lambda$ and the patrol pattern of length $c$ is repeated indefinitely. Thus, the expected time to detection is at least $1/\lambda$, which can be achieved only if the patroller allocates all their effort at this location all the time.
Consequently, we have
\[
\int_0^c \frac{1}{g(t)} dt = \lambda r c.
\]


Using Lemma~\ref{lem: am>=hm} and the preceding, we can find a lower bound for our objective function in \eqref{eq:obj} as follows:
\begin{equation}
\frac{\int_0^c g(t) dt}{c} \geq \frac{c}{\int_0^c \frac{1}{g(t)} dt} = \frac{c}{\lambda r c} = \frac{1}{\lambda r}.
\label{eq:one_LB}
\end{equation}
In other words,  $1 / (\lambda r)$ is a lower bound for the expected time to detect an attacker at this location, and this lower bound is achieved if $g(t)$ is constant for $t \in [0,c]$.
By taking $h(t) = r$ for $t \in [0,c]$, the instantaneous detection rate at the location stays at $r \lambda$ at all times, so $g(t) = 1/ (\lambda r)$ for $t \in [0,c]$, which achieves the lower bound in \eqref{eq:one_LB}.
Consequently, the optimal policy is to take $h(t) = r$ for $t \in [0,c]$, which completes the proof.
\end{proof}

\end{document}
